\chapter{前沿方向与展望}\label{ch:24-frontier}

\chapterintro{前面二十三章讲的都是“今天能直接拿来用”的方法。这一章换个视角：看看这个领域正在往哪走。我们不会罗列所有论文，而是挑出四条主线——语义与学习、多模态融合、长期地图维护、学习式导航——说清它们各自解决什么问题、现在到了什么程度、离工程落地还有多远。读完这一章，你应该能判断：哪些“新方法”值得投入时间，哪些还只是实验室里的玩具。}

\section{从几何 SLAM 到语义与学习}\label{sec:24-semantic}

\subsection{语义 SLAM}

传统 SLAM 输出的地图是“几何的”：点、面、栅格，不知道哪里是墙、哪里是门、哪里是人。语义 SLAM 在几何之外再估计每个点的类别或每个物体的实例，好处有三：
\begin{itemize}
  \item \textbf{动态物体剔除}：知道“这是人”，就可以把人的点从建图中剔除，而不是靠几何一致性猜；
  \item \textbf{更高层的导航}：可以下“去门口”这种指令，而不只是“去坐标 $(3,2)$”；
  \item \textbf{回环与重定位}：语义描述子对光照、季节变化更鲁棒。
\end{itemize}
代价是引入了分割网络，算力开销大、标注依赖强、分割错误会污染地图。目前工程上常见的是“离线语义标注 + 在线几何定位”的折中。

\subsection{学习式特征与描述子}

第 12 章讲的 Scan Context 是手工设计的描述子。学习式方法用神经网络直接学“两帧是否同一地点”：
\begin{itemize}
  \item 视觉侧：SuperPoint、SuperGlue、LoFTR 等把特征提取与匹配都交给网络，在弱纹理、大视角变化下明显优于 ORB；
  \item 激光侧：PointNet++、OverlapPredator、 diffusion-based 描述子等直接处理无序点云。
\end{itemize}
问题是：模型需要大量数据、推理耗时、跨场景泛化不确定。比赛场景相对固定，手工描述子往往够用；跨季节、跨场地的长期建图才更能体现学习式方法的优势。

\subsection{神经隐式表示}

NeRF（神经辐射场）和 3D Gaussian Splatting（3DGS）用神经网络或高斯集合表示场景，能渲染出照片级的新视角图像。把它们用于 SLAM 是一个活跃方向：用位姿 + 图像联合优化辐射场，同时得到高保真地图和更精确的位姿。
\begin{itemize}
  \item 优点：地图紧凑（网络权重或高斯参数）、可渲染、可做光度一致性约束；
  \item 缺点：训练/优化慢，实时性差，几何精度不一定优于点云，且对大规模场景显存压力大。
\end{itemize}
目前它更适合“离线高精地图重建”或“小场景实时”，离“比赛场上 10 Hz 定位”还有距离。

\subsection{端到端里程计}

DROID-SLAM 这类方法把特征提取、匹配、位姿优化放进一个可微的迭代网络，端到端训练。它在公开数据集上精度很高，但泛化到新传感器、新场景仍需微调，且算力需求远高于传统方法。对算法组而言，更现实的用法是“用学习式方法做前端的一部分（如特征、描述子、深度补全），后端仍用经典优化”。

\section{多模态融合}\label{sec:24-multimodal}

第 9 章讲过 LiDAR + IMU 的紧耦合。前沿方向是把更多模态塞进同一个估计器：
\begin{itemize}
  \item \textbf{LiDAR-视觉-惯性}：FAST-LIVO2\cite{yuan2024fastlivo2}、R3LIVE\cite{lin2022r3live} 把点到面残差与光度误差联合优化，在几何退化（长走廊、隧道）或视觉退化（黑暗、强光）时互相补位。
  \item \textbf{事件相机}：输出异步“亮度变化”事件，延迟极低、动态范围极大，适合高速运动与高动态场景；缺点是数据稀疏、传统视觉算法要重写。
  \item \textbf{4D 毫米波雷达}：能测距离、方位、俯仰、速度，雨雾烟尘下仍工作；缺点是点云稀疏、噪声大。
  \item \textbf{UWB 与 GNSS}：作为全局约束，与 LIO 融合（第 23 章 Z-Localization 已有实践）。
\end{itemize}

多模态融合的核心难点不是“把数据放一起”，而是\textbf{时间同步、外参标定、噪声建模、失效检测}。第 26 章讲的标定，在模态越多时越重要。

\section{长期建图与地图维护}\label{sec:24-lifelong}

比赛只需要一张场地地图，但真实应用（仓储、园区、家庭）需要机器人在几个月甚至几年里持续使用同一张地图。这就带来“长期建图”问题：
\begin{itemize}
  \item \textbf{多会话 SLAM}：多次建图的结果要能拼接、对齐到同一坐标系；
  \item \textbf{变化检测}：家具挪了、货架变了、墙刷了新漆，地图要能更新而不是被旧地图“带偏”；
  \item \textbf{地图压缩与版本管理}：点云地图动辄几个 GB，需要压缩、分层、增量加载；
  \item \textbf{定位鲁棒性}：场景变化后，描述子检索的召回率会下降，需要“变化感知”的重定位。
\end{itemize}
Z-Localization 的“先验地图包 + 描述子 + 可选 UWB”已经是一个简化版的多会话方案；工业界更完整的方案通常包含地图服务、版本管理、云端更新。

\section{学习式导航}\label{sec:24-learnnav}

第 18--21 章的规划与控制都是“经典方法”：图搜索、采样、优化、MPC。前沿方向是用学习来做决策：
\begin{itemize}
  \item \textbf{学习式局部规划}：用强化学习或模仿学习直接输出速度指令，在复杂动态场景中可能比手工规则更灵活；缺点是可解释性差、安全性难保证、训练成本高。
  \item \textbf{端到端导航}：从传感器输入直接到动作输出（如“视觉导航”），在特定场景（走廊跟随、人群穿行）表现不错，但泛化与安全仍是问题。
  \item \textbf{语义/任务导航}：输入“去厨房拿杯子”，系统自己拆解成探索、定位、抓取子任务，这是具身智能的核心问题，目前仍在实验室阶段。
  \item \textbf{多机器人协同}：多机共享地图、协同探索、任务分配，难点在通信受限下的一致性。
\end{itemize}

一个务实的判断：\textbf{短期内，学习式方法更可能作为经典规划器的“插件”（如学习式代价函数、学习式启发、学习式恢复行为），而不是整体替换。}经典方法提供安全保证与可解释性，学习提供灵活性与泛化。

\section{对 RoboMaster 的启示}\label{sec:24-rm}

结合比赛场景，给算法组三条建议：
\begin{enumerate}
  \item \textbf{先把经典方法吃透。} FAST-LIO2/Point-LIO、Cartographer、Nav2、Z-Localization 这些开源项目已经覆盖了比赛 90\% 的需求。把它们的参数、失效模式、源码吃透，比追新论文的收益高得多。
  \item \textbf{关注“可恢复性”与“可观测性”。} 比赛里最致命的不是精度差 1 cm，而是定位飞了之后起不来。Z-Localization 的状态机与前后端分离思路，比任何新算法都值得抄。
  \item \textbf{谨慎引入学习式方法。} 如果要用，优先用在“离线”或“非安全关键”环节：离线建图清理、描述子训练、仿真环境生成。在线定位与规划仍应以经典方法为主。
\end{enumerate}

\begin{summary}

\begin{enumerate}
  \item 语义 SLAM 在几何之外增加类别/实例信息，利于动态剔除、高层导航与重定位，但依赖分割网络。
  \item 学习式特征、描述子、神经隐式表示在特定场景优于手工方法，但实时性、泛化与算力仍是瓶颈。
  \item 多模态融合（LiDAR-视觉-惯性、事件相机、4D 毫米波）是提升鲁棒性的主要方向，核心难点是同步与标定。
  \item 长期建图需要多会话对齐、变化检测、地图压缩与版本管理。
  \item 学习式导航短期内更可能作为经典规划器的插件，而非整体替代。
  \item 对比赛而言，吃透经典开源方案、重视可恢复性，比追新更重要。
\end{enumerate}

\end{summary}

\begin{reading}

\begin{itemize}
  \item \cite{yuan2024fastlivo2}、\cite{lin2022r3live}：多模态紧耦合的两个代表，读完能理解“视觉如何补几何”。
  \item \cite{barfoot2017state}：第 11 章之后讨论了学习与估计的结合，适合理性看待“端到端”。
  \item 3D Gaussian Splatting 原始论文与开源实现：理解神经隐式表示的能力与局限。
  \item 近三年的 ICRA/IROS/RSS 综述：跟踪语义 SLAM、长期建图、学习式规划的最新进展。
\end{itemize}


\end{reading}


